\documentclass{article}
\usepackage{neurips_2026}

\usepackage[utf8]{inputenc}
\usepackage[T1]{fontenc}
\usepackage{hyperref}
\usepackage{url}
\usepackage{booktabs}
\usepackage{amsfonts}
\usepackage{nicefrac}
\usepackage{microtype}
\usepackage{xcolor}
\usepackage{graphicx}

\title{AutoResearch for Traffic Congestion Prediction:\\
Segment-Relative Features Break a Labeling-Driven F1 Ceiling}

\author{%
  Sherry Huang \\
  STAT 390 Capstone \\
  University of Illinois Urbana-Champaign \\
  \texttt{xiyuh5@illinois.edu}
}

\begin{document}
\maketitle

\begin{abstract}
We present an AutoResearch loop applied to 30-minute-ahead congestion prediction
on Chicago road segments using the Chicago Traffic Tracker dataset (786K observations,
Sept.\ 2023). The loop iterates over model changes, keeps improvements, and reverts
regressions — all evaluated against a fixed validation split. Over 30 experiments we
identify and partially fix a structural labeling problem: the global congestion threshold
(30th percentile of training speeds, $\approx$23\,mph) mislabels $\approx$16\% of
training rows because different road segments have very different normal speed ranges.
Adding four segment-relative features — per-segment mean speed, standard deviation,
z-score, and relative ratio — breaks an F1 ceiling at 0.658 that no amount of
hyperparameter tuning had been able to cross. Switching from Random Forest to LightGBM
with full feature subsampling regularization raises the final validation F1 to
\textbf{0.6806} (precision 0.635, recall 0.733), a \textbf{+21.4\% gain} over the
logistic regression baseline (F1\,=\,0.560). The result is stable across random seeds
(mean 0.674, std 0.005).
\end{abstract}

\section{Introduction}

Traffic congestion prediction is a well-studied problem, but in practice most deployed
systems rely on real-time sensor feeds or external map data. This project asks a
different question: \textit{how much can a model learn from speed history alone, without
any external data, using an AutoResearch loop that iterates without human guidance?}

The research contract was formalized as: predict whether a given Chicago road segment
will be congested exactly 30 minutes from now, using only historical speed readings,
time-of-day features, and segment-relative statistics, evaluated by validation F1 score
on a held-out time-based split.

The loop was constrained to modify only \texttt{src/model.py} (model definition and
feature selection), while \texttt{src/run.py} (data split, congestion threshold, metric
computation, logging) was frozen throughout. This constraint guaranteed that all
30 experiments were evaluated identically, making the experiment archive fully comparable.

\section{Data, Metric, and Baseline}

\paragraph{Dataset.} Chicago Traffic Tracker (City of Chicago Open Data), September 2023.
786,420 observations, 22 columns. Each row is one speed reading for one road segment at
one 10-minute timestamp. The dataset contains a \texttt{SEGMENT\_ID} column identifying
individual road segments, which proved crucial.

\paragraph{Label definition.} A segment is labeled \texttt{congested} (positive class)
if its speed falls below the 30th percentile of all training-set speeds ($\approx$23\,mph).
This threshold is computed from training data only. The class ratio is 3.49:1
(77.7\% non-congested).

\paragraph{Data split.} Deterministic, time-based:
train\,=\,70\%, validation\,=\,15\%, test\,=\,15\% (locked, never accessed during development).

\paragraph{Prediction target.} Congestion status 3 time steps (30 minutes) ahead,
computed by \texttt{shift(-3)} per segment.

\paragraph{Features.} The frozen \texttt{run.py} precomputes a superset of 17 features:
\texttt{SPEED}, lags 1--6, \texttt{rolling\_mean\_3}, \texttt{rolling\_std\_3},
\texttt{speed\_diff}, \texttt{HOUR}, \texttt{DAY\_OF\_WEEK}, \texttt{MONTH}, and
four segment-relative features added in Week 5 (\texttt{segment\_mean\_speed},
\texttt{segment\_std\_speed}, \texttt{speed\_vs\_seg\_mean}, \texttt{speed\_zscore}).

\paragraph{Metric.} Validation binary F1 score (\texttt{sklearn.metrics.f1\_score},
\texttt{average='binary'}).

\paragraph{Baseline.} Logistic regression with lags 1--3 and time features:
F1\,=\,0.5598, precision\,=\,0.755, recall\,=\,0.445 (exp\_001).

\section{AutoResearch Loop Design}

Each iteration of the loop follows: (1) read current \texttt{model.py};
(2) propose one change; (3) edit \texttt{model.py}; (4) run \texttt{python src/run.py
"description"}; (5) if F1 improves, commit; else revert via \texttt{git checkout src/model.py}
and log \texttt{status=discard}. All results are appended to \texttt{experiments/results.csv}.

The loop is \textit{greedy}: it keeps the first configuration that beats the current best
and builds from there. This matches gradient descent in spirit — no lookahead, no
backtracking beyond the single-step revert. The frozen evaluator is the key discipline
mechanism: changing the metric after seeing a good result is impossible by design.

\section{Experiments and Results}

\subsection{Progression by Phase}

Table~\ref{tab:phases} summarizes the five phases of experiments.

\begin{table}[h]
\caption{F1 progression by experimental phase (30 experiments total, 0 crashes).}
\label{tab:phases}
\centering
\begin{tabular}{llcc}
\toprule
Phase & Key change & Best F1 & Gain \\
\midrule
Baseline (exp\_001) & Logistic regression & 0.5598 & --- \\
Model selection (exp\_002--007) & RF + class balance + extended lags & 0.6566 & +0.097 \\
Hyperparameter grid (exp\_008--013) & Depth, trees, undersampling ratio & 0.6585 & +0.002 \\
Threshold sweep (exp\_014--018) & Decision threshold 0.35--0.50 & 0.6585 & +0.000 \\
Segment features (exp\_019--021) & speed\_zscore, speed\_vs\_seg\_mean, seg.\ mean & 0.6727 & +0.014 \\
Model comparison (exp\_022--025) & LightGBM vs RF vs XGBoost vs HGB & 0.6780 & +0.005 \\
LightGBM tuning (exp\_026--030) & num\_leaves, subsampling, n\_estimators & \textbf{0.6806} & +0.003 \\
\bottomrule
\end{tabular}
\end{table}

\subsection{The Ceiling Discovery}

Experiments 3--13 confirmed a hard ceiling at F1\,$\approx$\,0.658.
Changing \texttt{n\_estimators} from 50 to 200, \texttt{max\_depth} from 4 to 12,
and the undersampling ratio from 1:1 to 2:1 all moved F1 by less than 0.003.
The root cause: the global 23\,mph threshold mislabels $\approx$16.3\% of training rows.
A highway segment with a typical speed of 45\,mph is called ``congested'' at 22\,mph
even if that is its normal rush-hour pace. A side street at 20\,mph is called
``not congested'' even if that is unusually fast for it. No model can learn correctly
from systematically incorrect labels.

The fix was to add segment-relative features computed from training data only
(no leakage): \texttt{segment\_mean\_speed} and \texttt{segment\_std\_speed} via
\texttt{groupby("SEGMENT\_ID")["SPEED"].agg(...)}, plus derived features
\texttt{speed\_zscore} and \texttt{speed\_vs\_seg\_mean}. This gives the model
road-specific context: instead of asking ``is 22\,mph slow?'' it can ask
``is 22\,mph slow \textit{for this road}?'' F1 jumped to 0.6727 in one step.

\subsection{Model Comparison and Final Tuning}

With segment features in place, four model families were compared head-to-head
(same features, same undersampling, same threshold\,=\,0.40):

\begin{table}[h]
\caption{Model comparison on identical setup (17 features, 2:1 undersample, threshold=0.40).}
\label{tab:models}
\centering
\begin{tabular}{lccc}
\toprule
Model & F1 & Precision & Recall \\
\midrule
Random Forest (exp\_020) & 0.6727 & 0.644 & 0.704 \\
HistGradientBoosting (exp\_022) & 0.6553 & 0.540 & 0.834 \\
XGBoost (exp\_023) & 0.6712 & 0.625 & 0.724 \\
LightGBM 14 features (exp\_024) & 0.6736 & 0.627 & 0.727 \\
LightGBM 17 features (exp\_025) & 0.6780 & 0.638 & 0.724 \\
LightGBM tuned (exp\_030, \textbf{final}) & \textbf{0.6806} & \textbf{0.635} & \textbf{0.733} \\
\bottomrule
\end{tabular}
\end{table}

LightGBM consistently edged Random Forest. Adding the three previously unused features
(\texttt{rolling\_std\_3}, \texttt{MONTH}, \texttt{segment\_std\_speed}) gave +0.004.
Hyperparameter tuning — reducing \texttt{num\_leaves} from 63 to 31,
setting \texttt{min\_child\_samples=10}, and adding 20\% feature and row subsampling —
gave a further +0.003.

\subsection{Stability}

To confirm the result is not a one-off run, the final model was evaluated across
5 random seeds: mean F1\,=\,0.674, std\,=\,0.005, range [0.668, 0.682].
The fixed-seed (42) result (0.6806) lies at the upper end but within 1.5 standard
deviations of the mean. The result is stable.

\section{What Worked and What Did Not}

\paragraph{Worked.} The most impactful single step was adding segment-relative features
(+0.014 F1). LightGBM marginally outperformed Random Forest on this dataset (+0.003).
Undersampling the majority class 2:1 outperformed \texttt{class\_weight="balanced"}
in F1 while achieving better precision. The frozen evaluator made all comparisons valid.

\paragraph{Did not work.} HistGradientBoosting underperformed Random Forest in both
attempts (exp\_005 without segment features, exp\_022 with). Threshold tuning below
0.40 reliably raised recall at the cost of F1. Hyperparameter tuning on the RF
(six experiments across depth and tree count) moved F1 by under 0.003 total.

\paragraph{Fundamental limits.} The label is still globally defined — the congestion
threshold is the 30th percentile of all training speeds, not a segment-relative
percentile. Fixing this would require changing the frozen evaluator. The residual
F1 gap from a theoretical maximum is largely attributable to this irreducible label noise.

\section{Conclusion}

The AutoResearch loop surfaced a non-obvious data quality issue — a global labeling
threshold that mislabels segment-specific speed patterns — and fixed it within the
constraints of a frozen evaluator. The final model achieves F1\,=\,0.6806 on the
validation set, +21.4\% above baseline, with a stable result confirmed across multiple
seeds. The test set has not been evaluated and remains locked for final presentation.

The clearest lesson: for structured tabular data, feature engineering that encodes
domain knowledge (here, road-specific speed norms) is more impactful than model
selection or hyperparameter search. The loop arrived at this insight through systematic
elimination — every failed tuning experiment narrowed the search space and pointed
toward the underlying labeling problem.

\begin{ack}
This project was completed as part of STAT 390 Capstone at the University of Illinois
Urbana-Champaign. The AutoResearch loop was assisted by Claude (Anthropic) as the
agent executing experiments, with the research design, problem identification, and
scope decisions made in collaboration with the student. Code repository:
\url{https://github.com/SherryHuang0311/traffic-autoresearch}.
\end{ack}

\section*{References}

\small

[1] City of Chicago. (2023). \textit{Chicago Traffic Tracker -- Historical Congestion
Estimates by Segment}. Chicago Open Data Portal.
\url{https://data.cityofchicago.org/Transportation/Chicago-Traffic-Tracker-Historical-Congestion-Esti/77hq-huss}

[2] Ke, G., Meng, Q., Finley, T., Wang, T., Chen, W., Ma, W., Ye, Q., \& Liu, T.-Y. (2017).
LightGBM: A highly efficient gradient boosting decision tree.
\textit{Advances in Neural Information Processing Systems}, 30.

[3] Breiman, L. (2001). Random forests.
\textit{Machine Learning}, 45(1), 5--32.

[4] He, H., \& Garcia, E. A. (2009). Learning from imbalanced data.
\textit{IEEE Transactions on Knowledge and Data Engineering}, 21(9), 1263--1284.

\newpage
\section*{NeurIPS Paper Checklist}

\begin{enumerate}

\item {\bf Claims}
    \item[] Answer: \answerYes{}
    \item[] Justification: The abstract and introduction state the F1 improvement claim,
    the segment-relative feature contribution, and explicitly note the label limitation.

\item {\bf Limitations}
    \item[] Answer: \answerYes{}
    \item[] Justification: Section 5 explicitly discusses the global threshold limitation,
    the residual label noise, and the fact that fixing it requires changing the frozen evaluator.

\item {\bf Theory assumptions and proofs}
    \item[] Answer: \answerNA{}
    \item[] Justification: This is an empirical ML study; no theoretical results are claimed.

\item {\bf Experimental result reproducibility}
    \item[] Answer: \answerYes{}
    \item[] Justification: Full code is available at
    \url{https://github.com/SherryHuang0311/traffic-autoresearch}. All 30 experiments
    are logged in \texttt{experiments/results.csv} with descriptions and metrics.
    \texttt{random\_state=42} is fixed throughout. The data split is deterministic.

\item {\bf Open access to data and code}
    \item[] Answer: \answerYes{}
    \item[] Justification: Code is public on GitHub. Data is publicly available from
    the City of Chicago Open Data Portal (no license restrictions on government data).

\item {\bf Experimental setting/details}
    \item[] Answer: \answerYes{}
    \item[] Justification: Data split fractions, congestion threshold definition,
    feature computation, and hyperparameters are all specified in the paper and code.

\item {\bf Experiment statistical significance}
    \item[] Answer: \answerYes{}
    \item[] Justification: Stability across 5 random seeds is reported (mean, std, range)
    in Section 4.4. Error bars are not reported in tables as they span 30 non-repeated runs.

\item {\bf Experiments compute resources}
    \item[] Answer: \answerYes{}
    \item[] Justification: All experiments ran on a MacBook CPU (no GPU). Each run
    completed in 2.7--7.3 seconds. Total compute for 30 experiments: under 3 minutes.

\item {\bf Code of ethics}
    \item[] Answer: \answerYes{}
    \item[] Justification: This research uses publicly available government traffic data
    and poses no ethical risks.

\item {\bf Broader impacts}
    \item[] Answer: \answerYes{}
    \item[] Justification: Accurate congestion prediction can improve routing and reduce
    idle emissions. A potential negative use would be discriminatory routing that
    disadvantages certain neighborhoods; this is a risk for any traffic prediction system.

\item {\bf Safeguards}
    \item[] Answer: \answerNA{}
    \item[] Justification: The model predicts road congestion from public data;
    no high-risk model release is involved.

\item {\bf Licenses for existing assets}
    \item[] Answer: \answerYes{}
    \item[] Justification: Chicago Traffic Tracker data is released under the City of
    Chicago's open data license. scikit-learn (BSD), LightGBM (MIT), XGBoost (Apache 2.0).

\item {\bf New assets}
    \item[] Answer: \answerYes{}
    \item[] Justification: The experiment log, model code, and AutoResearch loop are
    released on GitHub with an MIT license.

\item {\bf Crowdsourcing and research with human subjects}
    \item[] Answer: \answerNA{}
    \item[] Justification: No human subjects involved.

\item {\bf Institutional review board (IRB) approvals}
    \item[] Answer: \answerNA{}
    \item[] Justification: No human subjects involved.

\item {\bf Declaration of LLM usage}
    \item[] Answer: \answerYes{}
    \item[] Justification: Claude (Anthropic) was used as the AutoResearch agent
    executing experiments and editing \texttt{model.py}. LLM usage is a core,
    documented part of the methodology described throughout the paper.

\end{enumerate}

\end{document}
